\documentclass[10pt]{article}
\usepackage[utf8]{inputenc}
\usepackage[T1]{fontenc}
\usepackage{amsmath}
\usepackage{amsfonts}
\usepackage{amssymb}
\usepackage[version=4]{mhchem}
\usepackage{stmaryrd}
\usepackage{bbold}

\begin{document}
$217$

L17: Primitive Roots for prime numbers\\
Theorem 5.9 (Legendre): Let $p$ be a prime number and $d \in \mathbb{Z}$ with $d>0$ and $d(p-1$. Then there are exactly $\phi(d)$ incongruent integers having onder $d$ modulo $p$.\\
Proof: Given a positive divisor al of $p-1$, let $f(d)$ be the number of integers between land $p-1$ inclusive having order $d$. We must show that $f(d)=\phi(d)$ for all such $d$. By Proposition 5.1 we know that each integer between I and p-1 inclusive has order dividing $p-1$. Thus

$$
\sum f(d)=p-1
$$

$$
d(p-1, d>0
$$

By Theorem 3.5, we have that

$$
\sum \phi(d)=p-1 .
$$

$d(p-1) d>0$


\begin{equation*}
\sum f(d)=\sum \phi(d) . \tag{1}
\end{equation*}


$d \mid p-1, d>0$

$$
d \mid p-1, d>0
$$

In view of (1), to show $f(d)=\phi(d)$ for all positive divisors $d$ of $p-1$, it suffices to show that $f(d) \leq \phi(d)$ for all such $d$. If $f(d)=0$, then $f(d)<\phi(d)$ is obvious. Assume that $f(d)>0$. Then there is an integer between 1 and $p-1$ inclusive, say a, such that ord $p^{a}=d$. The $d$ integers $a, a^{2}, \ldots, a^{d}$ are incongruent modulo $p$ since or $d_{m} a=d$ (otherwise $a^{i} \equiv a^{d}(\bmod p)$ for some $i>j$, and after applying the inverse of $a^{j}(\bmod p)$ we have $a^{i-j} \equiv 1(\bmod p)$ with $1 \leq i-j<d$, contradiction!) Since each integer $a, a^{2}, \ldots, a^{d}$ is a solution to the congruence $x^{d}-1 \equiv 0(\bmod p)$, Proposition 5.8 implies\\
that the incongrvent solutions modulo $P$ of the congruence $x^{d} \equiv 1(\bmod p)$ are given precisely by $a, a^{2}, \ldots, a^{d}$. Exactly $\phi(d)$ of these $d$ integers have order $d$ modulop by Corollary 5.5. Thus $f(d)=\phi(d)$ ( which certainly implies that $f(d) \leqslant \phi(d)$ ), and So the proof is complete.

Corollary 5.10: Let $p$ be a prime number. Then there are exactly $\phi(p-1)$ incongruent primitive roots modulo $p$. In particular, primitive roots exist for all prime numbers.\\
Proof: A primitive root modulo $p$ is an integer of order $\phi(p)=p-1$. There are exactly $\phi(p-1)$ such incongruent integers by theorem 5.9.

\begin{itemize}
  \item Theovem 5.9 cesserts the existence of integers\\
having certain orders. However, the Theovem (4) does not construct these integers.\\
Example: - Let $p=7$. Theovem 5.9 implies that there exist integers of orders $1,2,3$ and 6 . (these are the divisors of $\phi(\lambda)=6$ ). Furthermore,
\end{itemize}

\begin{center}
\begin{tabular}{c|c|c}
ovder & \# of residue classes & residue classes \\
1 & $\phi(1)=1$ & 1 \\
2 & $\phi(2)=1$ & 6 \\
3 & $\phi(3)=2$ & 2,4 \\
6 & $\phi(6)=2$ & 3,5 \\
\end{tabular}
\end{center}

Exercise: Construct the analogons table for $p=13$

Example: Find all incongruent integers having order 6 modulo 19 .

Theorem 5.9 tells us that there will be $\phi(b)=2$ such integers. We want to avoid an exhaustive Search through the integers from 1 to 18. Lets\\
first prove that 2 is a primitive root modulo modulo 19. By Proposition 5.1 we need only test $2^{a}$ where $a \mid \phi(19)=18$. Thus $a=1,2,3,6,9,18$. $2^{1} \equiv 2(\bmod 19), 2^{2} \equiv 4(\bmod 19), 2^{3} \equiv 8(\bmod 19)$, $2^{6} \equiv 7(\bmod 19)$ and $2^{9} \equiv 18(\bmod 19)$. Thus ord $192=18$. Proposition 5.3 tells us that the first 18 positive integral powers of 2 will generate all reduced residue classes modulo 19. Proposition 5.4 allows us to compute the orders of these integral powers of 2. We have

$$
\operatorname{ord}_{19}\left(2^{i}\right)=\frac{\operatorname{ord}_{19} 2}{\left(\operatorname{ord}_{19} 2, i\right)}=\frac{18}{(18, i)} .
$$

We want ord $19\left(2^{i}\right)=6$, which happens if and only if $(18, i)=3$. Thus $i=3$ or $i=15$. we have

$$
2^{3} \equiv 8(\bmod 19),
$$

and

$$
\begin{aligned}
2^{15} \equiv\left(2^{4}\right)^{3} 2^{3} \equiv(-3)^{3} 2^{3} & \equiv(-27)(8) \\
& \equiv(-8)(8) \\
& \equiv 12(\bmod 19) .
\end{aligned}
$$

Thus 8 and 12 are the desired incongrment integers having osder 6 modulo 19.

\begin{itemize}
  \item Consult Table 5.1 for the least primitive roots modulo $p$ for some small primes.\\
The frequency of 2 as the least positive primitive root modulop suggests the following conjecture.\\
Conjecture: There are infinitely many primes $p$ for which 2 is the least positive primitive root modulo $\rho$.
\end{itemize}

Conjecture (Artin): If $r$ is a non-square integer other than -1 , then there are infivitely many primes $p$ for which $r$ is a primitive root modulo p.

\begin{itemize}
  \item Artin's Conjecture is still open! Heath-Brown proved in 1986 that there are at most two integers $r$ for which Artin's conjecture fails!
\end{itemize}

\section*{The Primitive Root Theorem:}
\begin{itemize}
  \item Over the course of the next few lectures we will characterise the integers that have a primitive root.
  \item The next two Propositions will limit the integers that we need to consider.\\
Proposition 5.11: There are no primitive roots modulo $2^{n}$ where $n \in \mathbb{Z}$ and $n \geqslant 3$.\\
Proof: Note that a primitive root modulo $2^{n}$ must be odd and have order $\phi\left(2^{n}\right)=2^{n}-2^{n-1} =2^{n-1}$. Let $a$ be any odd integer. To prove that no primitive root modulo $2^{n}$ exists, it suffices to show that
\end{itemize}

$$
a^{2^{n-2}} \equiv 1\left(\bmod 2^{n}\right) .
$$

we use induction on $n$. If $n=3$, we must (8) show that $a^{2} \equiv 1(\bmod 8)$. We must have $a \equiv 1,3,5,7(\bmod 8)$, and in all cases we have that $a^{2} \equiv 1(\bmod 8)$. Assume that $k \geqslant 3$, and that the desired result holds for $n=k$ i.e. $a^{2^{k-2}} \equiv 1\left(\bmod 2^{k}\right)$. We must show the desired result holds for $n=k+1$. By the induction hypothesis we have $a^{2^{k-2}}-1=2^{k} b$ for some $b \in \mathbb{Z}$. Then $a^{2^{k-2}}=2^{k} b+1$, and squaring both sides of the equation gives:

$$
\begin{aligned}
a^{2^{k-1}} & =\left(2^{k} b\right)^{2}+2\left(2^{k} b\right)+1 \\
& \equiv 2^{k+1}\left(2^{k-1} b^{2}+b\right)+1 \\
& \equiv 1\left(\bmod 2^{k+1}\right)
\end{aligned}
$$

as required.

Proposition 5.12: There are no primitive (9) roots modulo $m n$ where $m, n \in \mathbb{Z}, m, n>2$, and $(m, n)=1$.\\
Proof: We will prove that the order of any integer modulo $m n$ is less than or equal to $\phi(m n) / 2$; since a primitive root modulo $m n$ must have order $\phi(m n)$, this suffices to prove the Proposition. Let $a \in \mathbb{Z}$ with $(a, m n)=1$.\\
Then $(a, m)=1$ and $(a, n)=1$. Put $d:=(\phi(m), \phi(n))$ and $k=[\phi(m), \phi(n)]$. Since $\phi(m)$ and $\phi(n)$ are both even we have $d \geqslant 2$. Thus

$$
k=\frac{\phi(m) \phi(n)}{d}=\frac{\phi(m n)}{d} \leqslant \frac{\phi(m n)}{2} .
$$

By Euler's Theorem, we have that $a^{\phi(m)} \equiv 1(\bmod m)$. Taking the $\frac{\phi(n)}{d}$ th power of both sides of the congunence gives

$$
a^{k}=a^{\phi(m) \phi(n) / d}=\left(a^{\phi(m)}\right)^{\phi(n) / d}
$$

$\left.\equiv\right|^{\phi(n) / d} \equiv \mid(\bmod m)$. Similarly,\\
$a^{k} \equiv 1(\bmod n)$. Since $(m, n)=1$, we have\\
that $a^{k} \equiv 1(\bmod m n)$.\\
Thus the order of any integer modulomn does not exceed $k \leq \frac{\phi(m n)}{2}$, as required.

Corollary 5.13: Let $n \in \mathbb{Z}$ with $n>0$. If\\
a primitive root modulo $n$ exists, then\\
$n$ must be equal to $1,2,4, p^{m}$ or $2 p^{m}$ where $p$ is an odd prime and $m$ a positive\\
integer.\\
Proof: Left as an exercise.


\end{document}